\documentclass[12pt]{article}

\usepackage{setspace}
\onehalfspace

%\usepackage[utf8]{inputenc}
%\usepackage{t1enc}
%\usepackage[T1]{fontenc}

\usepackage{amsmath}


%\DeclareMathOperator{\Exp}{Exp}
%\def\Exp{{\mathrm{Exp}}}

%\DeclareMathOperator{\Norm}{{\cal N}}
%\def\Norm{{\cal N}}

%\DeclareMathOperator{\Binom}{Binom}
%\def\Binom{{\mathrm{Binom}}}

%\DeclareMathOperator{\E}{E}
%\def\E{{\mathrm{E}}}

%\DeclareMathOperator{\D}{D}
%\def\D{{\mathrm{D}}}

%\renewcommand{\P}{\operatorname{P}}
%\def\P{{\mathrm{P}}}

%\DeclareMathOperator{\cov}{cov}
%\def\cov{{\mathrm{cov}}}


%\DeclareMathOperator{\corr}{corr}
%\def\corr{{\mathrm{corr}}}


%\renewcommand{\d}{\operatorname{d\! }} %integral \! a nyero
%\def\d{{\mathrm{d\! }}}


%\DeclareMathOperator{\R}{\mathbb{R}}
%\def\R{{\mathbb{R}}}

%\DeclareMathOperator{\Unif}{{\cal U}}
%\def\Unif{{\cal U}}

%\DeclareMathOperator{\Poi}{{Poisson}}
%\def\Poi{{\mathrm{Poisson}}}

%\def\ldots{{...\:}}


\usepackage{moodle}


\begin{document}

\begin{quiz}{gyakorló}

\begin{multi}[feedback={wiki}]{impfvv4}
Egy $\xi\sim \Unif(a,b)$ szórásnégyzete: $\frac{(a+b)^2}{2}$
\item igaz
\item* hamis
\end{multi}
\begin{multi}[feedback={Diszkrét vv-nak nincs sfv-e}]{sfv4}
Diszkrét valségi változók sűrűségfüggvénye lépcsős függvény.
\item igaz
\item* hamis
\end{multi}
\begin{multi}[feedback={Def./tulajdonság}]{dvv11}
Egy $\xi$ valségi változó az $x_1,\ldots ,x_n$ értékeket veheti fel. 
Ekkor a szórásnégyzete:
$$
\D^2(\xi)=\E(\xi^2)-\E(\xi)
$$
\item igaz
\item* hamis
\end{multi}
\begin{multi}[feedback={$\P(\xi=a)\le\P(a\le\xi<a+\frac{1}{n})=F(a+\frac{1}{n})-F(a) \overset{n\to\infty}{\to} 0$}]{eofv11}
Bármely $\xi$-re és $a$-ra $\P(\xi= a)=0$, feltéve hogy $\xi$ eloszlásfüggvénye
folytonos.
\item* igaz
\item hamis
\end{multi}
\begin{multi}[feedback={Additívitás.}]{val7}
Ha véges sok elemi esemény van, akkor bármely $C$-re $\P(C)=\sum_{\omega_k\in C} \P( \{\omega_k \})$.
\item* igaz
\item hamis
\end{multi}
\begin{multi}[feedback={de-Morgan azonosság}]{esem2}
Az $\overline{A+B}$ és $\overline{A} \cdot \overline{B}$ események kizárják egymást.
\item igaz
\item* hamis
\end{multi}
\begin{multi}[feedback={Def.}]{felt2}
$A$-nak a $B$-re vonatkozó feltételes valségét $\P(A|B)=\frac{\P(A\cdot B)}{\P(B)}$ módon értelmezzük, feltéve hogy $\P(B)>0$.
\item* igaz
\item hamis
\end{multi}
\begin{multi}[feedback={Def.}]{impdvv6}
Azt mondjuk, hogy $\xi\sim \Binom(n,p)$, ha 
$$
\P(\xi=k)=\binom{n}{k}p^k(1-p)\ \ \ (k=0,1\ldots n)
$$
valamilyen $n>0$ és $p\in[0,1]$ esetén.
\item igaz
\item* hamis
\end{multi}
\begin{multi}[feedback={Def.}]{dvv8}
Egy $\xi$ valségi változó az $x_1,\ldots, x_n$ értékeket veheti fel és $p_k=\P(\xi=x_k)$. 
Ekkor a második momentuma:
$$
\E(\xi^2)=\sum_{k=1}^n x_k^2 p_k
$$
\item* igaz
\item hamis
\end{multi}
\begin{multi}[feedback={$B=A\cup (B\setminus A)$+additívitás}]{val4}
A valószínűség monoton, azaz $\P(A)\le \P(B)$, ha $A\subseteq B$.
\item* igaz
\item hamis
\end{multi}
\begin{multi}[feedback={Def.}]{eofv1}
Egy $F:\R\to \R$ pontosan akkor eloszlásfüggvény, ha monoton nemcsökkenő, 
balról folytonos, $\lim_{x\to+\infty}F(x)=1$ és $\lim_{x\to-\infty}F(x)=0$
\item* igaz
\item hamis
\end{multi}
\begin{multi}[feedback={$A=B$?}]{val3}
A valószínűség monoton, azaz $\P(A)<\P(B)$, ha $A\subseteq B$.
\item igaz
\item* hamis
\end{multi}
\begin{multi}[feedback={Teljes valség tétele.}]{felt5}
Ha $A_1,\ldots,A_n$ teljes eseményrendszer és $\P(A_k)>0,\ k=1\ldots n$, akkor minden $B$-re:
$$
\P(B)=\sum_{k=1}^n \P(B|A_k)\P(A_k)
$$
\item* igaz
\item hamis
\end{multi}
\begin{multi}[feedback={Az adott feltételek mellet a várható érték egy súlyozott közép.}]{dvv5}
Egy $\xi$ valségi változó az $x_1,\ldots, x_n$ értékeket veheti fel. Ekkor:
$$
\min_k(x_k)\le \E(\xi) \le \max_k(x_k)
$$
\item* igaz
\item hamis
\end{multi}
\begin{multi}[feedback={Def.}]{impfvv8}
Ha $\xi\sim \Norm(0,1)$ akkor a sűrűségfüggvénye $f(x)=\frac{e^{-\frac{x^2}{2}}}{\sqrt{2\pi}}$.
\item* igaz
\item hamis
\end{multi}
\begin{multi}[feedback={$\xi=\xi_1+\ldots+\xi_n$, ahol $\xi_k$ Bernoulli és $\E(\xi_k)=p$+$\E$-additívitása}]{impdvv2}
Egy $\xi\sim \Binom(n,p)$ várható értéke: $n+p$
\item igaz
\item* hamis
\end{multi}
\begin{multi}[feedback={Def.}]{sfv11}
Ha $\xi$-nek $F$ illetve $f$ az eloszlásfüggvény illetve sűrűségfüggvénye, akkor bármely $a<b$-re:
$$
F(b)-F(a)=\int_{a}^{b}f(x)\d x
$$
\item* igaz
\item hamis
\end{multi}
\begin{multi}[feedback={Def.}]{flenvv5}
Ha $\xi$ és $\eta$ diszkrét valségi változóknak létezik a szórása és nemnulla, akkor a 
korrelációjuk
$$
\corr(\xi,\eta)=\E(\xi\eta)-\E(\xi)\E(\eta)
$$
\item igaz
\item* hamis
\end{multi}
\begin{multi}[feedback={Nókoment.}]{esem5}
Ha $\P(A)=1$ akkor $A$ bekövetkezhet.
\item* igaz
\item hamis
\end{multi}
\begin{multi}[feedback={Def.}]{sfv6}
Ha $\xi$ sűrűségfüggvénye $f$, akkor bármely $a<b$-re:
$$
\P(a<\xi<b)=\int_{a}^{b}f(x)\d x
$$
\item* igaz
\item hamis
\end{multi}
\begin{multi}[feedback={Tétel.}]{flenvv3}
Ha $\xi$ és $\eta$ valségi változók függetlenek, akkor
$$
\E(\xi \eta)=\E(\xi)\E(\eta)
$$
feltéve, hogy létezik a szórásuk.
\item* igaz
\item hamis
\end{multi}
\begin{multi}[feedback={Def.}]{sfv8}
Ha van $\xi$-nek $f$ sűrűségfüggvénye, akkor bármely $b$-re:
$$
\P(\xi<b)=f(b)
$$
\item igaz
\item* hamis
\end{multi}
\begin{multi}[feedback={Tulajdonság,tétel.}]{flenvv9}
Ha a $\xi$ és $\eta$ valségi változóknak létezik a korrelációja, akkor $-1\le \corr(\xi,\eta)\le 1$.
\item* igaz
\item hamis
\end{multi}
\begin{multi}[feedback={Def.}]{impfvv9}
Ha $\xi\sim \Exp(\lambda)$ akkor a sűrűségfüggvénye 
$$
f(x)=
\begin{cases}
\lambda e^{-\lambda x}& \ \ x>0\\
0 & \text{máskor} 
\end{cases}
$$
\item* igaz
\item hamis
\end{multi}
\begin{multi}[feedback={Def.}]{felt4}
$A_1,\ldots,A_n$ teljes eseményrendszer, ha $\sum A_k=\Omega$ és a tagok páronként függetlenek.
\item igaz
\item* hamis
\end{multi}
\begin{multi}[feedback={Def.}]{impfvv11}
Ha $\xi\sim \Unif(a,b)$ akkor a sűrűségfüggvénye 
$$
f(x)=
\begin{cases}
\frac{1}{b-a} & \ \ a<x<b \\
0 & \text{máskor} 
\end{cases}
$$
\item* igaz
\item hamis
\end{multi}
\begin{multi}[feedback={A balról folytonosság kell.}]{eofv3}
Egy $F:\R\to [0,1]$ pontosan akkor eloszlásfüggvény, ha monoton nemcsökkenő, 
$\lim_{x\to+\infty}F(x)=1$ és $\lim_{x\to-\infty}F(x)=0$
\item igaz
\item* hamis
\end{multi}
\begin{multi}[feedback={de-Morgan azonosság}]{esem4}
Az $\overline{A} + \overline{B}$ maga után vonja az $\overline{A\cdot B}$ eseményt.
\item* igaz
\item hamis
\end{multi}
\begin{multi}[feedback={Diszkrét v.v.-ra ez nem feltétlen igaz, ha $a$-t pozitív valséggel felveszi.}]{eofv8}
Bármely $\xi$-re és $a<b$-re $\P(a< \xi <b)=F_{\xi}(b)-F_{\xi}(a)$.
\item igaz
\item* hamis
\end{multi}
\begin{multi}[feedback={wiki}]{impfvv2}
Egy $\xi\sim \Unif(a,b)$ várható értéke: $\frac{b-a}{2}$
\item igaz
\item* hamis
\end{multi}
\begin{multi}[feedback={Def.}]{impdvv7}
Azt mondjuk, hogy $\xi\sim \Poi(\lambda)$, ha 
$$
\P(\xi=k)=e^{-\lambda}\frac{\lambda^k}{k!}\ \ \ (k=0,1,2\ldots )
$$ 
valamilyen $\lambda>0$ esetén.
\item* igaz
\item hamis
\end{multi}
\begin{multi}[feedback={Ez csak egyenletesre igaz.}]{dvv4}
Egy $\xi$ valségi változó az $x_1,\ldots, x_n$ értékeket veheti fel. Ekkor a várható értéke:
$$
\E(\xi)=\frac{\sum_{k=1}^n x_k}{n}
$$
\item igaz
\item* hamis
\end{multi}
\begin{multi}[feedback={Definíció.}]{esem7}
Az $A$ és $B$ események pontosan akkor függetlenek, ha $\P(AB)=\P(A)\P(B)$.
\item* igaz
\item hamis
\end{multi}
\begin{multi}[feedback={Def.}]{dvv9}
Egy $\xi$ valségi változó az $x_1,\ldots, x_n$ értékeket veheti fel és $p_k=\P(\xi=x_k)$. 
Ekkor a második momentuma:
$$
\E(\xi^2)=\sum_{k=1}^n x_k p_k^2
$$
\item igaz
\item* hamis
\end{multi}
\begin{multi}[feedback={Def.}]{flenvv1}
Azt mondjuk, hogy $\xi$ és $\eta$ diszkrét valségi változók függetlenek, ha
$$
\P(\xi=x,\eta=y)=\P(\xi=x)\P(\eta=y)
$$
bármely $x,y\in \R$-re teljesül.
\item* igaz
\item hamis
\end{multi}
\begin{multi}[feedback={Def.}]{dvv6}
Egy $\xi$ valségi változó az $x_1,\ldots, x_n$ értékeket veheti fel. Ekkor a szórása:
$$
\D(\xi)=\sum_{k=1}^n \P(\xi=x_k)(x_k-\E(\xi))^2
$$
\item igaz
\item* hamis
\end{multi}
\end{quiz}

\end{document}

